\newpage
\subsubsection{Justificación del proveedor de nube: AWS sobre Azure y GCP}
\label{subsec:justificacion_aws}

La selección de Amazon Web Services (AWS) como proveedor de
infraestructura en la nube para este sistema no es arbitraria. Se
fundamenta en un análisis comparativo de los tres principales
proveedores del mercado: AWS, Microsoft Azure y Google Cloud Platform
(GCP), evaluando sus servicios equivalentes, modelos de precios,
madurez del ecosistema y adecuación al contexto del proyecto.

\textbf{Posición en el mercado.} AWS es el proveedor de nube líder
a nivel mundial con una cuota de mercado del 31\% en 2025, seguido
por Azure con 25\% y GCP con 11\%, según datos de Synergy Research
Group \cite{aws_vs_azure_gcp_2025}. Esta posición de liderazgo se
traduce en un ecosistema más maduro, mayor disponibilidad de
documentación, comunidad de soporte más amplia y mayor cantidad de
integraciones con herramientas de terceros como Micronaut, Angular
y GitHub Actions.

\textbf{Comparativa de servicios equivalentes.} La
Tabla~\ref{tab:comparativa_nubes} presenta los servicios equivalentes
de cada proveedor para los componentes utilizados en el sistema.

\begin{longtable}{|>{\raggedright\arraybackslash}m{3cm}|>
{\raggedright\arraybackslash}m{3.5cm}|>{\raggedright\arraybackslash}m{3.5cm}|>
{\raggedright\arraybackslash}m{3.5cm}|}
\caption{Comparativa de servicios equivalentes entre AWS, Azure y GCP (2025).}
\label{tab:comparativa_nubes} \\
\hline
\textbf{Componente} & \textbf{AWS} & \textbf{Microsoft Azure} & \textbf{Google Cloud (GCP)} \\
\hline
\endfirsthead
\multicolumn{4}{c}{{\bfseries Tabla \thetable\ -- continuación}} \\
\hline
\textbf{Componente} & \textbf{AWS} & \textbf{Microsoft Azure} & \textbf{Google Cloud (GCP)} \\
\hline
\endhead

Funciones serverless & AWS Lambda & Azure Functions & Cloud Functions \\
\hline
API Gateway & Amazon API Gateway & Azure API Management & Cloud API Gateway / Apigee \\
\hline
Base de datos relacional & Amazon RDS & Azure SQL Database & Cloud SQL \\
\hline
Almacenamiento de objetos & Amazon S3 & Azure Blob Storage & Cloud Storage \\
\hline
CDN & Amazon CloudFront & Azure CDN & Cloud CDN \\
\hline
Autenticación & Amazon Cognito & Azure Active Directory B2C & Firebase Authentication \\
\hline
ML / Versionado de modelos & Amazon SageMaker & Azure Machine Learning & Vertex AI \\
\hline
CI/CD & GitHub Actions + AWS CodeDeploy & Azure DevOps & Cloud Build \\
\hline
Certificados SSL/TLS & AWS Certificate Manager & App Service Managed Certificates & Google-managed SSL \\
\hline
Monitoreo & Amazon CloudWatch & Azure Monitor & Cloud Monitoring \\
\hline
WAF & AWS WAF & Azure Web Application Firewall & Cloud Armor \\
\hline
\end{longtable}

\textbf{Comparativa de capa gratuita.} La
Tabla~\ref{tab:free_tier_comparativa} compara los límites gratuitos
de los servicios más relevantes para el sistema entre los tres
proveedores \cite{aws_vs_azure_gcp_2025}.

\begin{longtable}{|>{\raggedright\arraybackslash}m{3cm}|>
{\raggedright\arraybackslash}m{3.5cm}|>{\raggedright\arraybackslash}m{3.5cm}|>
{\raggedright\arraybackslash}m{3.5cm}|}
\caption{Comparativa de capa gratuita entre AWS, Azure y GCP (2025).}
\label{tab:free_tier_comparativa} \\
\hline
\textbf{Servicio} & \textbf{AWS} & \textbf{Azure} & \textbf{GCP} \\
\hline
\endfirsthead
\multicolumn{4}{c}{{\bfseries Tabla \thetable\ -- continuación}} \\
\hline
\textbf{Servicio} & \textbf{AWS} & \textbf{Azure} & \textbf{GCP} \\
\hline
\endhead

Funciones serverless & \textbf{1M invocaciones/mes siempre gratis} & 1M invocaciones/mes (12 meses) & 2M invocaciones/mes siempre gratis \\
\hline
Almacenamiento objetos & \textbf{5 GB siempre gratis} & 5 GB (12 meses) & 5 GB siempre gratis \\
\hline
Base de datos relacional & 750 hrs db.t2.micro (12 meses) & 250 GB Azure SQL (12 meses) & 1 instancia db-f1-micro siempre gratis \\
\hline
CDN & \textbf{1 TB transferencia siempre gratis} & 15 GB (12 meses) & 10 GB/mes siempre gratis \\
\hline
Autenticación & \textbf{50,000 usuarios activos/mes siempre gratis} & 50,000 MAU (siempre gratis) & 10,000 usuarios/mes (Firebase) \\
\hline
API Gateway & \textbf{1M llamadas REST/mes siempre gratis} & No tiene capa gratuita específica & 2M llamadas/mes (Cloud Endpoints) \\
\hline
Certificado SSL & \textbf{Gratis e ilimitado con servicios AWS} & Gratis solo en App Service & Gratis con servicios GCP \\
\hline
\end{longtable}

\textbf{Modelo de precios pay-as-you-go.} AWS fue pionero en el
modelo de pago por uso en la nube pública, y sigue siendo el más
granular de los tres proveedores. Lambda factura por milisegundo de
ejecución, S3 por GB almacenado y por número de solicitudes, y
CloudFront por GB transferido, lo que garantiza que el sistema pague
exactamente por lo que consume sin desperdiciar recursos. Esto
contrasta con modelos de otros proveedores donde algunos servicios
tienen cargos mínimos independientemente del uso \cite{northflank_comparison}.

\textbf{Razones concretas para elegir AWS sobre Azure y GCP.}

\begin{itemize}
    \item \textbf{Madurez del ecosistema serverless:} AWS Lambda fue
    el primer servicio de funciones serverless del mercado, lanzado
    en 2014, con dos años de ventaja sobre Azure Functions (2016) y
    Google Cloud Functions (2016). Esta madurez se traduce en mejor
    documentación, más ejemplos de integración con Micronaut y mayor
    estabilidad del servicio \cite{cloudplexo_serverless}.

    \item \textbf{Integración nativa con GitHub Actions:} El pipeline
    CI/CD del proyecto utiliza GitHub Actions, que tiene integración
    oficial y más desarrollada con AWS que con GCP. Azure DevOps es
    el competidor directo pero implica adoptar todo el ecosistema
    Microsoft.

    \item \textbf{Amazon SageMaker vs alternativas:} Para el
    versionado del modelo de clasificación, SageMaker Model Registry
    ofrece una integración más directa con S3 y Lambda que Vertex AI
    (GCP) o Azure Machine Learning, simplificando el flujo de
    publicación del modelo \texttt{.tflite}.

    \item \textbf{CloudFront vs Azure CDN vs Cloud CDN:} Amazon
    CloudFront ofrece 1 TB de transferencia gratuita de forma
    permanente, mientras que Azure CDN solo cubre 15 GB durante los
    primeros 12 meses y Cloud CDN ofrece 10 GB/mes. Para el frontend
    Angular del sistema, CloudFront representa la opción más económica
    a largo plazo \cite{aws_vs_azure_gcp_2025}.

    \item \textbf{Amazon Cognito vs alternativas:} Cognito ofrece
    50,000 usuarios activos mensuales de forma gratuita e indefinida,
    comparable a Azure AD B2C pero con una integración más directa
    con API Gateway y Lambda sin configuración adicional.

    \item \textbf{Reservaciones de RDS:} AWS ofrece Reserved Instances
    para RDS con hasta un 72\% de descuento sobre el precio bajo
    demanda, permitiendo planificar costos a largo plazo con certeza.
    GCP ofrece descuentos similares mediante Committed Use Discounts,
    pero la flexibilidad de instancia de AWS es mayor \cite{northflank_comparison}.

    \item \textbf{Comunidad y soporte académico:} AWS cuenta con el
    programa \textit{AWS Educate} y \textit{AWS Academy}, que
    proporcionan créditos y recursos educativos para proyectos
    académicos, alineado con el contexto de titulación de este
    prototipo.

    \item \textbf{Azure no es la opción óptima para este stack:}
    Azure es la mejor opción cuando se trabaja dentro del ecosistema
    Microsoft (Windows Server, SQL Server, .NET, Azure DevOps). El
    stack del presente sistema, basado en Java 2 1/Micronaut, Angular,
    PostgreSQL y GitHub Actions, no se beneficia de las ventajas
    específicas de Azure, lo que haría su adopción innecesariamente
    compleja sin ventaja económica o técnica.

    \item \textbf{GCP no es la opción óptima para este stack:} GCP
    destaca en cargas de trabajo de análisis de datos masivos con
    BigQuery, orquestación de contenedores con GKE y modelos de ML
    con Vertex AI. Para un sistema de sincronización móvil con backend
    REST y base de datos relacional, estas ventajas no son relevantes
    y su ecosistema serverless es menos maduro que el de AWS para
    integraciones con Micronaut y PostgreSQL.
\end{itemize}

En conclusión, AWS representa la opción más adecuada para este proyecto
por su madurez en el ecosistema serverless, la amplitud de su capa
gratuita, el modelo pay-as-you-go granular, la integración nativa con
las herramientas del stack seleccionado y su trayectoria como líder
del mercado de nube pública con más de 18 años de operación continua
\cite{aws_vs_azure_gcp_2025}.